\chapter{Discussion and Limitations}
\label{ch:discussion}

\section{Interpretation}

The results show that a standard 3D~U-Net, given the four co-registered MRI
modalities and trained with a combined Dice and cross-entropy objective, learns to
localise and segment glioblastoma sub-regions from multi-parametric MRI, reaching a
foreground mean Dice of 0.71. As is typical on BraTS, performance is not uniform
across sub-regions. The peritumoral edema achieves the highest Dice (0.77): it is
comparatively large and well contrasted on FLAIR. The enhancing tumour is also
segmented well (Dice 0.73) and in fact has the tightest boundary agreement (the
lowest HD95, 6.4~mm), reflecting the sharp contrast it shows on T1Gd. The
necrotic/non-enhancing core is the hardest sub-region (Dice 0.61): it is smaller,
more irregular and defined largely by the \emph{absence} of enhancement, which
makes its boundary intrinsically ambiguous. This ordering is consistent with the
qualitative example, where the core is the part most often under-segmented.

\section{Failure cases}
\label{sec:failure-cases}

Two failure modes are worth noting. First, in slices with little or no tumour the
model occasionally produces scattered false-positive voxels; this is characteristic
of under-training and of the class imbalance between background and tumour, and it
is reduced both by longer training and by the Dice term in the loss. Second, where
two tissue types have similar appearance --- for instance necrotic core against
non-enhancing tumour, or edema against normal white-matter hyperintensity --- the
model can mislabel one as the other. These are precisely the boundaries that also
cause disagreement between human raters. The worst validation case in
Figure~\ref{fig:cases} (Dice 0.22) is a clear instance of the second mode: the
model correctly localises the abnormal region via its edema but fails to resolve
the internal necrotic core and enhancing tumour, collapsing what should be three
distinct sub-regions into one.

\section{Comparison with the original 2023 project}

This work supersedes an earlier undergraduate project report on the same broad
topic, submitted in 2023. That report trained a 2D classifier on 253 curated MRI
images sorted only into ``tumour''/``no tumour'' folders, reporting an
image-level classification accuracy of 90.4\% together with a confusion matrix
whose reported counts were internally inconsistent, and without a clearly
documented dataset provenance, train/validation protocol, or a segmentation
evaluation metric beyond a handful of example masks. The present work differs
in ways that matter for scientific validity rather than for the headline
number alone: it trains and evaluates on the real, multi-institutional
BraTS~2020 dataset rather than a small curated image set; it performs full 3D
volumetric segmentation into four classes rather than binary image-level
classification; it reports standard, reproducible overlap and boundary metrics
(Dice, HD95) with a fixed, documented, seed-controlled train/validation split
(Section~\ref{sec:data}); and it is accompanied by a public, working software
artifact (Chapter~\ref{ch:conclusion}) rather than a static report alone. The
two projects' headline numbers are consequently not comparable to one another,
but the pipeline described here is markedly more rigorous, and every claim in
it can be independently re-run and checked.

\section{Limitations}

Several limitations should be stated plainly.

\begin{itemize}
  \item \textbf{Single dataset, no external test set.} The model is trained and
        evaluated only on BraTS~2020 HGG cases using an internal validation split.
        It has not been tested on data from other institutions or scanners, so its
        generalisation is unverified. A truly held-out or cross-institutional test
        set would be needed to make claims about clinical robustness.
  \item \textbf{Validation, not the official BraTS test protocol.} Results are
        reported on our own validation split, not through the official BraTS
        online evaluation portal or its held-out test set. Section~\ref{sec:literature-comparison}
        computes the standard composite regions for a like-for-like comparison
        with published numbers, but this remains an internal split rather than
        the official challenge evaluation.
  \item \textbf{Fixed centre-crop.} Volumes are centre-cropped to a $128^3$ patch
        for memory reasons. For most subjects this comfortably contains the brain,
        but a tumour lying far from the centre could in principle be partly
        clipped; a sliding-window or whole-volume inference scheme would remove
        this risk.
  \item \textbf{Restriction to HGG.} Focusing on glioblastoma keeps the framing
        accurate but reduces the training set to 293 subjects and means the model
        is not trained to handle lower-grade gliomas.
  \item \textbf{Compute budget.} Training on a single 8~GB GPU constrains the
        patch size, batch size and model capacity; larger patches, larger batches
        or transformer encoders were out of scope here.
\end{itemize}

These limitations are directions for improvement rather than flaws in the method,
and they inform the future work outlined next.
